En este capítulo se presentan los resultados obtenidos con las dos implementaciones desarrolladas: la versión \textit{offline}, que opera sobre registros de voltaje almacenados, y la versión \textit{online}, que ejecuta la optimización en tiempo real mediante módulos de \ac{rtxi}. Para cada caso se muestran tanto las dinámicas conseguidas como datos específicos del proceso de optimización bayesiana que permiten verificar su correcto funcionamiento. En ambas versiones se han utilizado 50 muestras iniciales de \ac{lhs} y 350 iteraciones de optimización bayesiana, lo que totaliza 400 evaluaciones. Considerando que el espacio de búsqueda comprende 9 parámetros en el caso \textit{offline} y 10 en el \textit{online}, este presupuesto resulta adecuado: las 50 muestras iniciales proporcionan una cobertura suficiente del espacio mediante \ac{lhs}, y las 350 iteraciones restantes permiten que el \ac{gp} converja hacia regiones óptimas sin resultar excesivas. Los rangos de corriente objetivo son distintos para fase y antifase y para la versión \textit{offline} y \textit{online}, y han sido determinados experimentalmente para cada caso con el fin de producir la dinámica buscada sin generar artefactos por exceso de intensidad.

Cae destacar que, como muchos de los resultados que se presentarán para ambas versiones \textit{offline} y \textit{online} (tanto en fase como en antifase) son muy similares, muchas de las explicaciones se omiten para evitar redundancia. Además, debido a esto o a la carencia de interés analítico de algunas de estas gráficas (en específico las de las importancias de los parámetros, como se comentará), se han desplazado las mismas al Apéndice~\ref{CAP:FIGURASADICIONALES}.


\section{Resultados de la implementación \textit{offline}\label{SEC:RESULTADOS_OFFLINE}}

La versión \textit{offline} se ha evaluado tanto para acoplamiento en antifase (inhibitorio) como en fase (excitatorio), optimizando simultáneamente $I_{\mathrm{fast}}$ e $I_{\mathrm{slow}}$ en dirección unidireccional (se inyecta su suma). El paso temporal del CSV es de 0{,}1\,ms (dado que se extrajo con una frecuencia de muestreo de 10\,kHz), con un tiempo de estabilización de 4000\,ms y tiempos de evaluación y observación (para el escalado) de 7000\,ms, abarcando unos 5 periodos de \textit{burst}. Al tratarse de un registro estático que se puede leer directamente y sin las esperas de adquisición de las neuronas vivas, estos tiempos prolongados no suponen un coste temporal significativo.

La frecuencia de corte del filtro de separación frecuencial se ha establecido en $0{,}024$\,kHz, valor verificado experimentalmente para separar adecuadamente los \textit{spikes} de la onda lenta en este registro. La corrección de \textit{drift} (Sección~\ref{SUBSUBSEC:ESCALADO}) se ha desactivado para no introducir deformaciones en la señal estática del CSV.

Dado el bajo tiempo de optimización en ambos casos (ahora se ven), se han podido ejecutar 20 repeticiones independientes y promediar los resultados, aumentando la significación estadística de las observaciones.


\subsection{Antifase\label{SUBSEC:OFFLINE_ANTIFASE}}

Para esta optimización, el rango de corriente total objetivo se ha fijado entre $-1{,}4$ y $-0{,}4$\,nA adimensionales (al estar en la escala de la neurona Hindmarsh-Rose). El tiempo de optimización medio ha sido de $29{,}57$\,s, con una desviación estándar de $9{,}27$\,s, lo cual es bajo, como ya se comentó.

La Figura~\ref{FIG:OFFLINE_ACC_MAX_ANTI} muestra la evolución de la puntuación máxima acumulada promediada sobre las 20 ejecuciones. La curva comienza entre $0{,}55$ y $0{,}60$, crece rápidamente durante las muestras iniciales de \ac{lhs} hasta estabilizarse entre $0{,}70$ y $0{,}75$, y experimenta un nuevo incremento pronunciado poco después de concluir las 50 muestras iniciales (en torno a la evaluación 60), estabilizándose entre $0{,}85$ y $0{,}90$. Esta transición refleja la dinámica de la optimización bayesiana: el estancamiento inicial ocurre porque el objetivo del \ac{lhs} no es maximizar la puntuación, sino aportar una cobertura inicial del espacio de búsqueda. Una vez que el \ac{gp} dispone de datos suficientes para ser informativo, la función \ac{ei} guía la exploración de forma efectiva, propiciando el nuevo ascenso hasta el estancamiento final de la convergencia. Aunque la curva se estabiliza de forma evidente, conserva una ligera tendencia ascendente, lo que sugiere la posibilidad de obtener mejoras marginales con iteraciones adicionales. Esta mejora continua demuestra un desempeño adecuado del algoritmo con el presupuesto de evaluaciones: no resulta excesivo, puesto que el modelo sigue refinando la solución, ni insuficiente, dada la clara convergencia alcanzada. Finalmente, la variabilidad se mantiene aproximadamente constante y moderada a lo largo del proceso, demostrando una consistencia robusta del algoritmo entre las distintas repeticiones.

\begin{figure}[Promedio de la puntuación máxima acumulada por evaluación \textit{offline} en antifase]{FIG:OFFLINE_ACC_MAX_ANTI}{Evolución del promedio de la puntuación máxima acumulada por evaluación a lo largo de 20 ejecuciones independientes para el acoplamiento en antifase utilizando la implementación \textit{offline}. La región sombreada indica la desviación estándar (como si fuesen barras de error verticales). La línea discontinua vertical separa las muestras iniciales de \ac{lhs} de las iteraciones guiadas por la optimización bayesiana.}
    \image{0.9\textwidth}{}{bo_history_offline_syn_i_antiphase_acc_max}
\end{figure}

La Figura~\ref{FIG:OFFLINE_AVG_SCORES_ANTI} desglosa la puntuación de la corriente promediada en sus componentes de forma y rango. Ambas componentes muestran un perfil fluctuante en torno a unos valores concretos, sin una clara tendencia general al alza o a la baja dentro de cada fase (muestras iniciales e iteraciones). Aunque promediar múltiples ejecuciones enmascara la progresión individual de cada ruta de muestreo y optimización (aplana la curva), se puede observar el peso relativo de cada componente. La puntuación de forma se sitúa consistentemente por encima de la de rango (poco por debajo de $0{,}8$ frente a $0{,}6$), especialmente tras las muestras iniciales (ambas estaban en torno a $0{,}6$, aunque la de rango ligeramente por debajo) tras las que mejora la puntuación por lo descrito en el anterior párrafo. Esto indica que el algoritmo encuentra combinaciones de parámetros que reproducen bien la morfología de $V_{\mathrm{pre}}$, mientras que ajustar el rango de la corriente resulta más difícil, probablemente debido al fenómeno del \textit{parameter sloppiness} descrito en la Sección~\ref{SUBSUBSEC:PARAMETER_SLOPPINESS}, donde se comenta que este afecta en gran medida a la amplitud y desplazamiento de la corriente. Las barras de error de la puntuación de rango son ligeramente menos estables (aunque de media son parecidas), reflejando menor consistencia entre ejecuciones, quizá también debida a este fenómeno (porque el algoritmo no identifica bien en qué dirección avanzar) o a que el ajuste de la normalización de la puntuación de corriente descrito en la Sección~\ref{SUBSUBSEC:PUNTUACION_RANGO} pudiese ser algo más preciso. Sin embargo, esta normalización cumple su propósito al mantener ambas componentes en rangos bastante parecidos (aunque no iguales) en media y variabilidad, evitando que una domine sobre la otra. La puntuación total se sitúa entre ambas, como corresponde a su cálculo como media aritmética (ligeramente por debajo de$0{,}7$ en las iteraciones).

\begin{figure}[Promedio de la puntuación y sus componentes por evaluación \textit{offline} en antifase]{FIG:OFFLINE_AVG_SCORES_ANTI}{Promedio de la puntuación total y de sus componentes (puntuación de forma y puntuación del rango de corriente) por evaluación a lo largo de 20 ejecuciones para el acoplamiento en antifase en la implementación \textit{offline}. La región sombreada indica la desviación estándar de cada curva (como si fuesen barras de error verticales). La línea discontinua vertical separa las muestras iniciales de \ac{lhs} de las iteraciones guiadas por la optimización bayesiana.}
    \image{0.9\textwidth}{}{bo_history_offline_syn_i_antiphase_avg_scores}
\end{figure}

La Figura~\ref{FIG:OFFLINE_SCORES_ANTI} muestra las puntuaciones individuales de una ejecución representativa en formato de dispersión. En general aplica lo mismo que en la versión promediada, pero se pueden ver algunos detalles adicionales. La mayoría de los puntos de ambas componentes se concentran entre $0{,}5$ y $0{,}85$, con la puntuación de rango tendiendo a la parte baja de este intervalo y la de forma a la parte alta. Los \textit{outliers} por debajo de esta zona, evidencian los episodios de exploración inducidos por la función de adquisición \ac{ei}. La ausencia de \textit{outliers} por arriba sugiere que las puntuaciones altas ya están próximas al límite alcanzable para esta configuración. Debido a esta mayoría de puntos cercanos al límite superior encontrado sugiere que el algoritmo es predominantemente explotador, afirmación coherente con el valor bajo de \textit{jitter} ($\xi = 0{,}003$) (Sección~\ref{SUBSUBSEC:DIFERENCIAS_OFFLINE_ONLINE}).

\begin{figure}[Puntuación y sus componentes por evaluación \textit{offline} en antifase]{FIG:OFFLINE_SCORES_ANTI}{Valores individuales de la puntuación total y de sus componentes (puntuación de forma y puntuación del rango de corriente) obtenidos en cada evaluación para una ejecución representativa del acoplamiento en antifase en la implementación \textit{offline}. La línea discontinua vertical separa las muestras iniciales de \ac{lhs} de las iteraciones guiadas por la optimización bayesiana.}
    \image{0.9\textwidth}{}{bo_history_offline_syn_i_antiphase_scores}
\end{figure}

La Figura~\ref{FIG:OFFLINE_ARD_ANTI} muestra la importancia promediada de cada parámetro, cuantificada como el inverso de la escala de longitud (\textit{lengthscale}) del kernel \ac{se}-\ac{ard} del \ac{gp} tras la optimización (lo que el algoritmo ha aprendido sobre cada parámetro individualmente). Una escala de longitud corta (importancia alta) indica que el objetivo varía rápidamente al modificar ese parámetro con respecto al rango de variación del mismo (porque están normalizados con respecto a este), identificándolo como una dirección \textit{stiff}; una escala larga (importancia baja) señala una dirección \textit{sloppy} en la que el objetivo es insensible (Sección~\ref{SUBSEC:OPTIMIZACION_BAYESIANA}). Sin embargo, extraer conclusiones absolutas a partir de esta jerarquía resultaría especulativo debido a dos sesgos importantes: la fuerte dependencia de las \textit{lengthscales} respecto a los rangos de búsqueda predefinidos (ya que se normalizan en base a los mismos), y la distorsión intrínseca al comparar parámetros en escala lineal con aquellos en escala logarítmica (que miden sensibilidades a incrementos absolutos frente a multiplicativos, respectivamente). Por consiguiente, no se realiza un análisis de estos resultados, pero se mantiene la figura por su interés para visualizar que el algoritmo es capaz de aprender una jerarquía en la importancia de cada parámetro y cuál es esta jerarquía.

\begin{figure}[Promedio de la importancia en la puntuación de los parámetros \textit{offline} en antifase]{FIG:OFFLINE_ARD_ANTI}{Importancia relativa promediada de cada parámetro sobre la puntuación objetivo a lo largo de 20 ejecuciones para la optimización \textit{offline} en antifase. Se calcula como el inverso de las escalas de longitud (\textit{lengthscales}) del kernel \ac{ard} tras finalizar la optimización. Las barras más altas indican parámetros frente a los cuales la función objetivo es más sensible. Las barras de error representan la desviación estándar. Cada \textit{lengthscale} está normalizado en base a su rango de búsqueda.}
    \image{0.9\textwidth}{}{bo_history_offline_syn_i_antiphase_ard_lss}
\end{figure}

En cuanto a los valores concretos de los parámetros obtenidos, no se detallan en el texto por estar en la escala de la neurona Hindmarsh-Rose y no ser varios de ellos directamente interpretables; no obstante, el promedio y desviación estándar de cada parámetro en estas ejecuciones se pueden consultar en el repositorio de datos de las figuras. A pesar de esto, se puede ver que se cumplen relativamente lo que se explicó teóricamente en la Sección~\ref{SUBSEC:SINAPSIS_QUIMICA}: $E_{\mathrm{syn}}$ se sitúa establemente lejos del rango de $V_{\mathrm{post}}$ (para disminuir la influencia de $V_{\mathrm{post}}$ respecto a la de $V_{\mathrm{pre}}$) y por debajo (su rango entero de variación en antifase está por debajo para causar inhibición); $g_{\mathrm{fast}}$ es menor que $g_{\mathrm{slow}}$, acorde con que los \textit{spikes} del registro de la neurona \ac{pd} usado son de menor amplitud relativa que la onda lenta; $k_1$ es menor que $k_2$ (y por tanto $R$ es menor que $1$), lógico para que \(m_{\mathrm{slow}}\) permanezca elevado a lo largo del \textit{burst} y decaiga entre \textit{bursts}; $V_{\mathrm{fast}}$ se establece en la parte alta del rango de $V_{\mathrm{pre}}$ para centrar la sigmoide en los \textit{spikes}; y $s_{\mathrm{fast}}$ tiende a valores altos dentro del rango para comprimir la onda lenta en la parte inferior de la sigmoide rápida. En el caso de $s_{\mathrm{slow}}$, ya que su rango está bastante ajustado, los valores que tome dentro de este no son demasiado informativos. Sin embargo, $V_{\mathrm{slow}}$ no parece ajustarse a los valores esperados, presentando una variabilidad bastante grande dentro del rango en el que se definió (que es rango más lógico), y tendiendo a valores altos dentro del rango de \(V_{\mathrm{pre}}\). Esto que puede deberse a que, como $s_{\mathrm{slow}}$ será pequeño (porque su rango es de valores pequeños), aunque se desplace el umbral y, por tanto, el rango de entradas de la sigmoide a un extremo de la misma, sigue sin deformarse. Es decir, con este rango de $s_{\mathrm{slow}}$,  \(V_{\mathrm{pre}}\) actúa de manera \textit{sloppy}, ya sea de forma individual al no afectar apenas a la puntuación, o porque otros parámetros compensen fácilmente los cambios de \(V_{\mathrm{pre}}\). Esto es el fenómeno de \textit{parameter sloppiness} descrito, en el que varias combinaciones de parámetros pueden dar resultados similares. Además, es posible que también haya una posible compensación entre $k_2$ y $k_1$, ya que ambos presentan una variabilidad bastante grande con respecto a sus valores típicos. El resto de parámetros tienen una variabilidad moderada o baja con respecto a estos valores típicos, lo que da mayor confianza en lo explicado en este párrafo sobre los mismos.

Finalmente, la Figura~\ref{FIG:OFFLINE_DYN_ANTI} presenta la dinámica resultante de una simulación con parámetros obtenidos de una optimización representativa (la misma ejecución que la de la gráfica de las componentes de la puntuación). La corriente total reproduce la forma esperada: la corriente lenta elimina los \textit{spikes} de $V_{\mathrm{pre}}$ (filtrados mediante $k_1$, $k_2$ y/o la compresión de los mismos en la parte alta de la sigmoide), mientras que la rápida los reproduce comprimiendo la onda lenta en su parte baja. Su suma reconstituye la morfología completa de $V_{\mathrm{pre}}$. En cuanto al rango de la corriente total, se aproxima considerablemente al objetivo; lo que sugiere que alcanzar el límite exacto podría no ser posible o requeriría un compromiso con la forma de la corriente, aunque el usado es aceptable porque se queda cerca. La inhibición provoca que $V_{\mathrm{post}}$ se hunda ligeramente durante el disparo presináptico, impidiendo que dispare a la vez que la neurona presináptica y causando la antifase. Sin embargo, tanto en la componente rápida como en la lenta se aprecian los \textit{spikes} hacia arriba de $V_{\mathrm{post}}$. Estas alteraciones son inevitables ya que $V_{\mathrm{post}}$ actúa directamente sobre la corriente, sin sigmoides o ecuaciones diferenciales de por medio, aunque han sido reducidas eficazmente empleando un valor grande de $E_{\mathrm{syn}}$.

\begin{figure}[Potenciales de membrana y corriente sináptica en la dinámica resultante \textit{offline} en antifase]{FIG:OFFLINE_DYN_ANTI}{Resultado de la simulación \textit{offline} en antifase con una combinación de parámetros obtenida de una optimización representativa. La gráfica superior muestra los potenciales de membrana de la neurona presináptica (del registro CSV de la neurona \ac{pd}) y de la postsináptica (modelo Hindmarsh-Rose). Las gráficas inferiores desglosan la corriente sináptica en su componente total, componente rápida y componente lenta, ilustrando cómo inducen el comportamiento en antifase. Tanto los potenciales como las corrientes operan en la escala de la neurona Hindmarsh-Rose ya que el escalado se hace a la neurona del registro (presináptica)}
    \image{0.9\textwidth}{}{offline_syn_i_antiphase}
\end{figure}


\subsection{Fase\label{SUBSEC:OFFLINE_FASE}}

Para este caso, el rango de corriente total objetivo se ha establecido entre $0{,}2$ y $0{,}8$\,nA adimensionales. El tiempo de optimización medio ha sido de $32{,}2$\,s, con una desviación estándar de $6{,}77$\,s.

La evolución de la puntuación máxima acumulada para el caso de fase es prácticamente idéntica a la del caso de antifase (Figura~\ref{FIG:OFFLINE_ACC_MAX_FASE_AP}), aplicando las mismas explicaciones.

En cuanto a las puntuaciones promediadas, el comportamiento es muy similar también (Figura~\ref{FIG:OFFLINE_AVG_SCORES_FASE_AP}), aplicando explicaciones similares también, aunque hay una variabilidad ligeramente mayor en la componente de rango para la antifase.

Con las puntuaciones individuales sucede lo mismo, pero ahora con algún \textit{outlier} más lejano en la zona baja y más puntos en torno a $0{,}5$ para la puntuación de rango (Figura~\ref{FIG:OFFLINE_SCORES_FASE_AP}).

En cuanto a la importancia de los parámetros (Figura~\ref{FIG:OFFLINE_ARD_FASE_AP}), los resultados no son muy distantes a los de antifase, lo que es lógico porque el único parámetro que debe cambiar si se busca fase o antifase es $E_{\mathrm{syn}}$ ya que este cambia el signo de la corriente, pero la forma o magnitud, controlada por los otros parámetros, no debe cambiar. Dado que su rango tiene la misma amplitud dependiente del rango de $V_{\mathrm{post}}$ para fase y antifase, y está igual de alejado del rango de $V_{\mathrm{post}}$ (por encima o por debajo en función de si es fase o antifase), su importancia en la puntuación debería ser parecida (para la puntuación de rango solo si el rango de corriente buscado para antifase está bien definido con respecto al de fase).

El hecho de que haya cambios entre estas gráficas puede deberse a la aleatoriedad propia de cada ejecución así como a que el rango de corriente buscado en fase, como se dijo, no esté definido a la perfección con respecto al de antifase (aunque sí que lo está bastante bien porque hay pocas variaciones).

En cuanto a los parámetros obtenidos, $E_{\mathrm{syn}}$ se sitúa ahora por encima del rango de $V_{\mathrm{post}}$ también de manera alejada, por los mismos motivos, pero ahora logrando un comportamiento excitatorio en lugar del inhibitorio. El resto de parámetros tienen un comportamiento similar, lo que es lógico, porque, como se explicó, $E_{\mathrm{syn}}$ es el único parámetro que debe cambiar para pasar de antifase a fase.

La Figura~\ref{FIG:OFFLINE_DYN_FASE} muestra la dinámica resultante en fase también para una simulación con parámetros obtenidos de una optimización representativa (la misma ejecución que la de la gráfica de las componentes de la puntuación). También aquí se puede ver que las corrientes del resultado obtenido son adecuadas tanto en forma como en rango, pero ahora siendo ambas simétricas respecto al eje horizontal con respecto a las de fase (ahora hay en forma una correlación negativa con respecto a $V_{\mathrm{pre}}$) debido a la inversión de signo causada por el cambio del valor de $E_{\mathrm{syn}}$. De esta forma, cuando la neurona presináptica dispara, se inyecta más corriente a la postsináptica, propiciando la fase (que ambas disparen a la vez). El resto de explicaciones son análogas a las del caso de antifase, salvo el efecto que causa esta corriente en la neurona postsináptica que se acaba de explicar.

\begin{figure}[Potenciales de membrana y corriente sináptica en la dinámica resultante \textit{offline} en fase]{FIG:OFFLINE_DYN_FASE}{Resultado de la simulación \textit{offline} en fase con una combinación de parámetros obtenida de una optimización representativa. La gráfica superior muestra los potenciales de membrana de la neurona presináptica (del registro CSV de la neurona \ac{pd}) y de la postsináptica (modelo Hindmarsh-Rose). Las gráficas inferiores desglosan la corriente sináptica en su componente total, componente rápida y componente lenta, ilustrando cómo inducen el comportamiento en fase. Tanto los potenciales como las corrientes operan en la escala de la neurona Hindmarsh-Rose ya que el escalado se hace a la neurona del registro (presináptica)}
    \image{0.9\textwidth}{}{offline_syn_i_phase}
\end{figure}


\section{Resultados de la implementación \textit{online}\label{SEC:RESULTADOS_ONLINE}}

La versión \textit{online} se ha ejecutado según la arquitectura de módulos de \ac{rtxi} descrita en la Sección~\ref{SUBSUBSEC:ARQUITECTURA_MODULOS} (Figura~\ref{FIG:MODULOS_MODELO}) para dos neuronas modelo Hindmarsh-Rose interconectadas bidireccionalmente. En este caso se ha optimizado $I_{\mathrm{fast}}$ en un sentido (de la neurona 1 a la 2) e $I_{\mathrm{slow}}$ en el otro (de la neurona 2 a la 1), en lugar de optimizar ambas componentes en ambos sentidos. Esta decisión reduce el número de parámetros a optimizar simultáneamente (que de otro modo sería elevado, aunque factible con un presupuesto de evaluaciones suficiente), y a la vez permite verificar la bidireccionalidad y el correcto funcionamiento de ambas componentes en direcciones distintas. Las neuronas Hindmarsh-Rose se han configurado con un periodo de \textit{burst} de aproximadamente $1{,}6$\,s, valor que es normal en algunas neuronas biológicas. Se ha empleado un tiempo de estabilización de $400$\,ms y un tiempo de evaluación de $1800$\,ms, ligeramente superior a un periodo, lo que garantiza al menos un \textit{burst} completo para valorar la similitud de la corriente con su componente frecuencial de $V_{\mathrm{pre}}$ correspondiente, siendo suficiente por los motivos expuestos en la Sección~\ref{SUBSEC:FUNCION_EVALUACION}.

La frecuencia de corte del filtro se ha establecido en $0{,}002$\,kHz, valor que se ha comprobado experimentalmente que permite separar correctamente los \textit{spikes} y la onda lenta de la neurona Hindmarsh-Rose con este periodo de \textit{burst}. Cabe destacar que se tenía un periodo de \ac{rtxi} (cada vez que se ejecuta la función de tiempo real) de $1$\,ms ($1$\,kHz de frecuencia de muestreo). Con esto, se puede verificar la relación $f_c / f_s$ (frecuencia de corte entre frecuencia de muestreo) es parecida a la empleada en la versión \textit{offline}, lo que garantiza una separación frecuencial parecida. Sin embargo, no es exactamente igual debido a que al escalar temporalmente la neurona Hindmarsh-Rose sus spikes pueden ensancharse o estrecharse de manera que duren distinto a los de una neurona viva (como la \ac{pd} usada en la versión \textit{offline}).

Ahora, dado que cada evaluación requiere esperar el tiempo real de estabilización y adquisición, las ejecuciones son considerablemente más largas que en la versión \textit{offline}. Por esta magnitud en los tiempos y porque el funcionamiento del algoritmo ya fue analizado con mayor significación estadística (20 repeticiones) en la versión \textit{offline} y aquí el comportamiento debería ser parecido, no se han realizado promedios de varias ejecuciones, mostrando solo una. Los valores concretos de los parámetros obtenidos no se comentan al carecer de rigurosidad estadística por ser una sola ejecución y también por haber sido ya explicados en la Sección~\ref{SUBSEC:OFFLINE_ANTIFASE}.


\subsection{Antifase\label{SUBSEC:ONLINE_ANTIFASE}}

Para esta optimización, los rangos de corriente objetivo se han fijado entre $-1{,}0$ y $0{,}0$\,nA adimensionales para $I_{\mathrm{fast}}$ y entre $-1{,}4$ y $-0{,}4$ para $I_{\mathrm{slow}}$. El tiempo de optimización ha sido de $915{,}46$\,s (poco más de 15 minutos), un tiempo aceptable para un experimento hipotético con una neurona viva.

La Figura~\ref{FIG:ONLINE_ACC_MAX_ANTI} muestra la evolución de la puntuación máxima acumulada. La curva comienza poco por encima de $0{,}74$ y asciende durante las muestras iniciales de \ac{lhs} hasta llegar a $0{,}83$. Este valor se mantiene hasta aproximadamente la evaluación 125, donde se encuentra una solución superior a $0{,}9$ (mejor que los vistos hasta ahora en los promedios de la versión \textit{offline}) que se conserva prácticamente hasta el final, con una mejora marginal casi inapreciable en una iteración cercana al término. Las explicaciones sobre la dinámica de convergencia expuestas en la Sección~\ref{SUBSEC:OFFLINE_ANTIFASE} aplican de igual modo al ser un comportamiento parecido, aunque con menor significación estadística al tratarse de una única ejecución.

\begin{figure}[Puntuación máxima acumulada por evaluación \textit{online} en antifase]{FIG:ONLINE_ACC_MAX_ANTI}{Evolución de la puntuación máxima acumulada por evaluación durante la ejecución en tiempo real (\textit{online}) para el acoplamiento en antifase. La línea discontinua vertical separa las muestras iniciales de \ac{lhs} de las iteraciones guiadas por la optimización bayesiana.}
    \image{0.9\textwidth}{}{bo_history_rtxi_syn_model_ifast_islow_antiphase_acc_max}
\end{figure}

La Figura~\ref{FIG:ONLINE_SCORES_ANTI} muestra las puntuaciones individuales por evaluación. El comportamiento es análogo al descrito para la versión \textit{offline} (Sección~\ref{SUBSEC:OFFLINE_ANTIFASE}), pero con la mayoría de los valores concentrados en un rango más alto ($0{,}6$ a $0{,}9$), con menos \textit{outliers} por debajo, y con ambas componentes de la puntuación más segregadas dentro de este rango: los puntos de forma se agrupan cerca de $0{,}9$ y los de rango cerca de $0{,}6$, quedando la puntuación total poco por debajo de $0{,}8$. Estas tres observaciones se explican por el valor más bajo de \textit{jitter} ($\xi = 0{,}001$) en la versión \textit{online} (Sección~\ref{SUBSUBSEC:DIFERENCIAS_OFFLINE_ONLINE}), que hace al algoritmo más explotador: al buscar soluciones próximas a la mejor encontrada, las puntuaciones se mantienen en un rango más alto, con menor dispersión y con las componentes más estables en sus respectivas regiones.

\begin{figure}[Puntuación y sus componentes por evaluación \textit{online} en antifase]{FIG:ONLINE_SCORES_ANTI}{Valores individuales de la puntuación total y de sus componentes (forma y rango) obtenidos en cada evaluación durante la ejecución en tiempo real (\textit{online}) para el acoplamiento en antifase. La línea discontinua vertical separa las muestras iniciales de \ac{lhs} de las iteraciones guiadas por la optimización bayesiana.}
    \image{0.9\textwidth}{}{bo_history_rtxi_syn_model_ifast_islow_antiphase_scores}
\end{figure}

La Figura~\ref{FIG:ONLINE_ARD_ANTI_AP} (también en el Apéndice~\ref{CAP:FIGURASADICIONALES}) presenta la importancia de los parámetros para esta ejecución. Por las mismas razones expuestas en la Sección~\ref{SUBSEC:OFFLINE_ANTIFASE}, no se analiza en profundidad, pero se incluye para evidenciar que el algoritmo aprende una jerarquía de importancias. Hay parámetros cuya importancia es prácticamente nula, y los dos parámetros más importantes coinciden con los que también lo eran en la versión \textit{offline}. Esto razonable al tratarse del mismo algoritmo, pese a que ahora cada componente de la corriente opere en una dirección distinta y una de las neuronas también sea diferente (ahora es un modelo cuando antes era un registro de una viva).

Finalmente, la Figura~\ref{FIG:ONLINE_DYN_ANTI} muestra la dinámica resultante. La separación frecuencial de cada señal Hindmarsh-Rose se ha realizado correctamente: $I_{\mathrm{fast}}$ (de la neurona 1 a la 2) reproduce los \textit{spikes} suprimiendo la onda lenta, e $I_{\mathrm{slow}}$ (de la neurona 2 a la 1) reproduce la envolvente lenta filtrando los \textit{spikes}, aunque en su traza se aprecia inevitablemente la forma de $V_{\mathrm{post}}$ por la misma razón que en la versión \textit{offline} (la influencia directa de $V_{\mathrm{post}}$ en la corriente), pero estando también mitigado mediante un $E_{\mathrm{syn}}$ alto. La forma de cada corriente se asemeja a la componente frecuencial correspondiente de su $V_{\mathrm{pre}}$ respectivo (que no es el mismo para ambas al operar cada sinapsis en una dirección distinta). En cuanto al efecto sobre los voltajes, la corriente rápida inyectada provoca ligeros \textit{spikes} invertidos (hacia abajo) en la neurona receptora, mientras que la corriente lenta produce un hundimiento suave (en a otra neurona, ya que son sinapsis en sentidos opuestos); ambos efectos contribuyen a impedir que las neuronas disparen simultáneamente, consiguiendo la antifase. En cuanto a los rangos de las corrientes, aquí parece que se ha más complicado reproducir los puestos como objetivo que en la versión \textit{offline} (están más lejos), quizá debido, como se dijo, a una definición inexacta de los mismos o a falta de tiempo para la optimización. Sin embargo, esto no ha impedido que se pueda conseguir la dinámica buscada.

\begin{figure}[Potenciales de membrana y corrientes sinápticas en la dinámica resultante \textit{online} en antifase]{FIG:ONLINE_DYN_ANTI}{Resultado de la ejecución en tiempo real (\textit{online}) en antifase con una combinación de parámetros obtenida de una optimización. La gráfica superior muestra los potenciales de membrana de ambas neuronas modelo Hindmarsh-Rose. Las gráficas inferiores muestran la componente rápida (de la neurona 1 a la 2 o \verb|1->2|) y la componente lenta (de la neurona 2 a la 1 o \verb|2->1|) de las corrientes sinápticas que inducen el comportamiento en antifase. Todos los valores operan en la escala adimensional de las neuronas Hindmarsh-Rose.}
    \image{0.9\textwidth}{}{rtxi_syn_model_ifast_islow_antiphase}
\end{figure}


\subsection{Fase\label{SUBSEC:ONLINE_FASE}}

En esta configuración, los rangos de corriente objetivo se han fijado entre $0{,}0$ y $0{,}6$\,nA adimensionales para $I_{\mathrm{fast}}$ y entre $0{,}2$ y $0{,}8$ para $I_{\mathrm{slow}}$. El tiempo de optimización ha sido de $934{,}14$\,s (casi 16 minutos), de nuevo un tiempo aceptable para interactuar con una neurona biológica.

La Figura~\ref{FIG:ONLINE_ACC_MAX_FASE} presenta la evolución de la puntuación máxima acumulada para el caso de fase. La curva comienza en torno a $0{,}65$ y asciende rápidamente de forma escalonada durante las muestras iniciales hasta poco menos de $0{,}85$, una muy buena puntuación para estar en el muestreo inicial. A partir de ahí, la puntuación sube progresivamente hasta poco antes de la evaluación 300, donde se encuentra la mejor solución, que es cercana a $0{,}9$ por debajo. El hecho de que esta se localice relativamente lejos del final y de que las muestras iniciales de \ac{lhs} ya proporcionen una puntuación elevada se debe a la aleatoriedad inherente a una sola ejecución, pero el comportamiento general no difiere sustancialmente del descrito en la versión \textit{offline} para antifase. Por tanto, aplican las mismas explicaciones que en las secciones previas, con la misma salvedad de la menor significación estadística.

\begin{figure}[Puntuación máxima acumulada por evaluación \textit{online} en fase]{FIG:ONLINE_ACC_MAX_FASE}{Evolución de la puntuación máxima acumulada por evaluación durante la ejecución en tiempo real (\textit{online}) para el acoplamiento en fase. La línea discontinua vertical separa las muestras iniciales de \ac{lhs} de las iteraciones guiadas por la optimización bayesiana.}
    \image{0.9\textwidth}{}{bo_history_rtxi_syn_model_ifast_islow_phase_acc_max}
\end{figure}

Las puntuaciones individuales (Figura~\ref{FIG:ONLINE_SCORES_FASE_AP}, Apéndice~\ref{CAP:FIGURASADICIONALES}) muestran un comportamiento análogo al de antifase, aunque con los puntos de la puntuación de forma ligeramente menos concentrados en la zona alta cercana al $0{,}9$ y con algún \textit{outlier} adicional de menor puntuación, si bien estos siguen siendo muy escasos. Aplican, por tanto, explicaciones similares que en antifase.

La importancia de los parámetros (Figura~\ref{FIG:ONLINE_ARD_FASE_AP}, Apéndice~\ref{CAP:FIGURASADICIONALES}) presenta igualmente importancias prácticamente nulas en algunos parámetros, y, aunque las barras difieren entre fase y antifase, las tres de mayor magnitud mantienen una proporción similar en ambos casos, cosa que probablemente se explique igual que se explicó en la Sección~\ref{SUBSEC:OFFLINE_FASE}. Además, los dos parámetros más importantes son coincidentes con los de la versión \textit{offline}, igual que pasaba con antifase, por el mismo motivo (se trata del mismo algoritmo).

La Figura~\ref{FIG:ONLINE_DYN_FASE} muestra la dinámica resultante en fase. De manera análoga a la antifase, la separación frecuencial y la forma de las corrientes son adecuadas. En $I_{\mathrm{slow}}$, la influencia de $V_{\mathrm{post}}$ es especialmente visible, dificultando la apreciación de la onda lenta, aunque esta se encuentra correctamente reproducida. En este caso, las corrientes provocan que cuando una neurona dispara, la otra reciba excitación, propiciando el disparo simultáneo y consiguiendo la fase. Sin embargo, ocurre lo mismo en cuanto a los rangos de las corrientes: no se han encontrado de manera exacta, probablemente debido a los mismos motivos.

\begin{figure}[Potenciales de membrana y corrientes sinápticas en la dinámica resultante \textit{online} en fase]{FIG:ONLINE_DYN_FASE}{Resultado de la ejecución en tiempo real (\textit{online}) en fase con una combinación de parámetros obtenida de una optimización. La gráfica superior muestra los potenciales de membrana de ambas neuronas modelo Hindmarsh-Rose. Las gráficas inferiores muestran la componente rápida (de la neurona 1 a la 2 o \verb|1->2|) y la componente lenta (de la neurona 2 a la 1 o \verb|2->1|) de las corrientes sinápticas que inducen el comportamiento en fase. Todos los valores operan en la escala adimensional de las neuronas Hindmarsh-Rose.}
    \image{0.9\textwidth}{}{rtxi_syn_model_ifast_islow_phase}
\end{figure}
